% Part: lambda-calculus
% Chapter: introduction
% Section: computable-lambda

\documentclass[../../../include/open-logic-section]{subfiles}

\begin{document}

\olfileid{lam}{int}{lrp}
\olsection{قابلِ حساب فنکشن \usetoken{S}{lambda definable} نیں}

\begin{thm}
\ollabel{thm:computable-lambda}
ہر جزوی قابلِ حساب فنکشن !!{lambda definable} اے۔
\end{thm}

\begin{proof}
سانوں دکھانا اے کہ ہر جزوی قابلِ حساب فنکشن~$f$ کسے
لامبڈا ٹرم~$F$ نال !!{lambda defined} اے۔ کلینی دے
نارمل روپ دے قضیے دے مطابق ایہ دکھانا کافی اے کہ ہر
بنیادی بازگشتی فنکشن کسے لامبڈا ٹرم نال !!{lambda defined}
اے۔ فیر ایہ دکھاؤنا اے کہ !!{lambda definable} فنکشن موزوں
ترکیباں تے بغیر مقررہ حد دی تلاش ہِٹھ بند نیں۔ہر بنیادی بازگشتی
فنکشن نوں کسے لامبڈا ٹرم نال !!{lambda defined} دکھاؤن
لئی ایہ دکھانا کافی اے کہ ابتدائی فنکشن !!{lambda definable}
نیں،تے اوہ جزوی فنکشن جیہڑے !!{lambda definable} نیں
ترکیب، بنیادی بازگشت تے بغیر مقررہ حد دی تلاش ہِٹھ بند نیں۔
\end{proof}

باقی ثبوت پڑھنا سہل بناؤن لئی اسیں زیادہ رائج لکھت ورتاں گے۔
مثلاً،$M(x, y, z)$ نوں $Mxyz$ دی تھاں لکھاں گے۔ایہ لکھت
بھاویں اک اشارہ دیندی اے،پر یاد رکھو کہ بلا نوع بندی لامبڈا
کیلکولس دے ٹرم دیاں دلیلاں دی کوئی مقرر گنتی نہیں ہوندی۔ ایس لئی اوسے ٹرم~$M$ لئی $M(x,y)$ تے
$M(x,y,z,w)$ لکھنا وی اکّو طرح جائز اے۔پر ایہ لکھت
دسدی اے کہ اسیں $M$ نوں تِن متغیراں والا فنکشن سمجھ کے ورت رہے آں۔
ایس نال تعریفاں دا مطلب وی صاف رہندا اے۔
ایس طرح اسی ``$M$ نوں $M(x,y,z) = \dots$ نال تعریف کرو''
آکھاں گے،بجائے ایس دے کہ ``$M$ نوں
$M = \lambd[x][\lambd[y][\lambd[z][\dots]]]$ نال تعریف کرو''۔

\end{document}
